% ID: SEL-D0016-BC-E05
% ESTADO: validada. S017. Revisión matemática, didáctica y visual.
\textbf{a) Probabilidad de error.}\par\nopagebreak[4]
Sean $P_i$ el suceso de que la persona $i$ realizó el cobro y $E$ el error. La tercera persona hizo el 25\,\%. La probabilidad total es
\[
\boxed{P(E)=0.30(0.01)+0.45(0.03)+0.25(0.02)=0.0215.}
\]

\textbf{b) Segunda persona entre los cobros correctos.}\par\nopagebreak[4]
El complemento tiene probabilidad $P(E^c)=0.9785$. Por la regla del producto, $P(P_2\cap E^c)=0.45(0.97)=0.4365$. La condicionada pedida es
\[
\boxed{P(P_2\mid E^c)=\frac{0.4365}{0.9785}=\frac{873}{1957}\approx0.4461.}
\]
El denominador también se obtiene sumando los cobros correctos de las tres personas.
